\documentclass[conference]{IEEEtran}
\IEEEoverridecommandlockouts

\usepackage{cite}
\usepackage{amsmath,amssymb,amsfonts}
\usepackage{algorithmic}
\usepackage{algorithm}
\usepackage{graphicx}
\usepackage{textcomp}
\usepackage{xcolor}
\usepackage{booktabs}
\usepackage{multirow}
\usepackage[colorlinks=true, linkcolor=black, citecolor=black, urlcolor=black]{hyperref}

\def\BibTeX{{\rm B\kern-.05em{\sc i\kern-.025em b}\kern-.08em
    T\kern-.1667em\lower.7ex\hbox{E}\kern-.125emX}}

\begin{document}

\title{Cache-Aided NOMA for 6G Networks: Deep Reinforcement Learning for\\
CIC-Aware Proactive Cache Placement}

\author{
\IEEEauthorblockN{Akash Mandal}
\IEEEauthorblockA{\textit{Dept. of Information Technology}\\
\textit{Techno Main Salt Lake}\\
Kolkata, India\\
13000222028@tmsl.ac.in}
\and
\IEEEauthorblockN{Arpita Guchhait}
\IEEEauthorblockA{\textit{Dept. of Information Technology}\\
\textit{Techno Main Salt Lake}\\
Kolkata, India\\
13000222031@tmsl.ac.in}
\and
\IEEEauthorblockN{Manojit Pal}
\IEEEauthorblockA{\textit{Dept. of Information Technology}\\
\textit{Techno Main Salt Lake}\\
Kolkata, India\\
13000222036@tmsl.ac.in}
\and
\IEEEauthorblockN{Vaibhav Anand}
\IEEEauthorblockA{\textit{Dept. of Information Technology}\\
\textit{Techno Main Salt Lake}\\
Kolkata, India\\
13000222037@tmsl.ac.in}
\and
\IEEEauthorblockN{Dr. Biswajit Ghosh}
\IEEEauthorblockA{\textit{Dept. of Information Technology}\\
\textit{Techno Main Salt Lake}\\
Kolkata, India \\
biswajit.ghosh@tmsl.ac.in}
\and
\IEEEauthorblockN{Mr. Abdul Aziz}
\IEEEauthorblockA{\textit{Dept. of Information Technology}\\
\textit{Techno Main Salt Lake}\\
Kolkata, India \\
abdul.aziz@tmsl.ac.in}
}

\maketitle

%=============================================================
% ABSTRACT
%=============================================================
\begin{abstract}
This paper proposes a Deep Reinforcement Learning (DRL) based proactive caching framework for Cache-aided Non-Orthogonal Multiple Access (NOMA) downlink systems targeting 6G networks. While caching and NOMA are promising complementary technologies, traditional reactive caching policies fail to optimize for Cache-aided Interference Cancellation (CIC) opportunities. The proposed scheme employs a Double Dueling Deep Q-Network (DQN) agent with a novel binary action space to make online cache placement decisions. By incorporating a CIC-aware reward function, the agent is explicitly incentivized to proactively cache files that enable interference suppression for paired NOMA users. Simulated under a realistic mixed fading channel model with imperfect Successive Interference Cancellation (SIC), results over 5 million training steps and 100 Monte Carlo evaluation runs demonstrate exceptional performance. The proposed DQN scheme achieves a cache hit rate of 70.38\%, a CIC benefit rate of 3.78 (a 5.28x improvement over the best dynamic baseline), and a sum-rate of 18.17 bps/Hz at an SNR of 30 dB. Furthermore, the proposed policy reduces system outage probability by 14x compared to the no-caching baseline, performing within 0.54\% of the oracle-optimal Top-K policy without relying on prior knowledge of the content popularity distribution.
\end{abstract}

\begin{IEEEkeywords}
cache-aided NOMA, deep reinforcement learning, 6G, proactive caching, CIC, Double DQN, interference cancellation, edge caching
\end{IEEEkeywords}


%=============================================================
% I. INTRODUCTION
%=============================================================
\section{Introduction}

Next-generation (6G) wireless networks are expected to meet unprecedented demands for ultra-low latency, massive connectivity, and high spectral efficiency \cite{b_ding_survey}. Two complementary technologies---Non-Orthogonal Multiple Access (NOMA) and edge caching---have independently demonstrated significant gains in spectral efficiency and backhaul reduction, respectively. NOMA enables multiple users to share the same time-frequency resources via power-domain multiplexing and Successive Interference Cancellation (SIC) \cite{b_dai_noma, b_ding_survey}, while edge caching reduces access latency and backhaul load by proactively placing popular content near the users \cite{b_maddah_ali, b_golrezaei}.

The joint paradigm---\textit{Cache-Aided NOMA}---has emerged as a highly promising architecture for 6G, where cached content can further enable Cache-aided Interference Cancellation (CIC). When a file requested by a paired NOMA user is locally available in the cache, the multi-user interference term from that user can be perfectly reconstructed and canceled without requiring error-prone SIC decoding. This fundamentally reduces outage probability and improves system-level fairness \cite{b_survey_bepari, b_cao_joint}.

However, the full potential of CIC depends critically on the \textit{quality} of the cached content---specifically, whether the cached files are the ones likely to be requested by paired NOMA users in the near future. Static and reactive policies (e.g., Top-K, Least Recently Used (LRU), Least Frequently Used (LFU)) do not account for this cross-layer NOMA coupling, resulting in severely suboptimal CIC exploitation. Deep Reinforcement Learning (DRL) offers an appealing solution: an intelligent agent can learn, through continuous environmental interaction, to dynamically cache files that simultaneously maximize the cache hit rate and CIC opportunities.

\subsection*{Contributions}
The primary contributions of this paper are summarized as follows:
\begin{itemize}
    \item We design a novel NOMA-aware state representation (14-dimensional) for the DRL agent that jointly encodes content popularity, cache occupancy, continuous channel statistics, and recent NOMA pairing history.
    \item We present a binary action space redesign that reduces the traditional DQN eviction action space from $\mathcal{O}(N_{\text{files}})$ to $\mathcal{O}(1)$, making DRL remarkably tractable for large, real-world content catalogs.
    \item We propose a specialized CIC-incentivized immediate reward function that explicitly rewards cache placement decisions enabling interference cancellation for paired NOMA users.
    \item We conduct a comprehensive simulation study under realistic mixed fading channels with imperfect SIC, comparing the DQN scheme against four baselines (Top-K Zipf, LRU, LFU, Random) and a No-Cache reference over a full signal-to-noise ratio (SNR) sweep from $-10$ to $30$ dB.
\end{itemize}


%=============================================================
% II. RELATED WORK
%=============================================================
\section{Related Work}

Edge caching for wireless systems has been extensively studied since early empirical observations demonstrated that content popularity follows Zipf-like distributions \cite{b_breslau_zipf}. Fundamental limits of caching---formalized in the coded-caching framework---proved that coordinated placement enables multicast opportunities and global caching gains \cite{b_maddah_ali}. Architectures such as Femtocaching and device-to-device (D2D) caching have subsequently been proposed to exploit spatial locality to reduce core network load \cite{b_golrezaei, b_ji_d2d}.

The convergence of edge caching and NOMA has recently been analyzed from stochastic geometry \cite{b_zhang_stochastic} and joint optimization perspectives \cite{b_cao_joint, b_guo_edge}. A comprehensive recent survey \cite{b_survey_bepari} identified CIC as a key unexplored dimension in cache-aided NOMA. While cached content can eliminate an interferer's signal at the receiver without SIC, this cross-layer benefit has not been fully exploited in conventional machine-learning-based placement policies.

Machine-learning approaches---ranging from statistical predictors to deep-RL methods---have been increasingly proposed for adaptive cache storage \cite{b_chen_ml, b_elghariani_pred}. Li \textit{et al.} \cite{b_li_drl} studied DRL-based dynamic caching in NOMA networks but did not incorporate a CIC-aware reward feedback mechanism or the binary action-space abstraction necessary for scalable catalog deployments. Cooperative and UAV-relay assisted NOMA designs have further been studied for cell-edge coverage \cite{b_bepari_partial, b_liu_uav}. Our work explicitly differs by formulating the first CIC-incentivized DRL proactive caching policy with direct NOMA-layer state feedback mapping.


%=============================================================
% III. SYSTEM MODEL
%=============================================================
\section{System Model}

\subsection{Network Topology}
We consider a single-cell downlink system with one Base Station (BS) and $K = 200$ users uniformly distributed within a cell of radius $R = 500$ m. The BS is equipped with a local cache capable of storing $C = 200$ files (representing 10\% of the entire catalog). A content catalog of $N = 2000$ unique files is available, with instantaneous request probabilities adhering to a Zipf distribution:

\begin{equation}
p_n = \frac{n^{-\alpha}}{\sum_{i=1}^{N} i^{-\alpha}}, \quad n = 1, \ldots, N
\label{eq:zipf}
\end{equation}

where $\alpha = 1.0$ dictates the skewness parameter of content popularity.

\subsection{Channel Model}
Each down-link user $k$ is characterized by a channel gain $h_k$, modeled as a product of distance-dependent path loss and small-scale fading. Path loss is formulated as $d_k^{-\beta}$ with a path loss exponent $\beta = 3.5$. To capture the heterogeneous propagation environments expected in 6G deployments, small-scale fading is modeled via a mixed fading distribution: Rayleigh fading with probability $(1 - p_{\text{LoS}})$ and Rician fading featuring a K-factor of $K_{\text{dB}} = 10$ dB with a line-of-sight probability $p_{\text{LoS}} = 0.4$.

\subsection{NOMA Transmission}
Users are organized into NOMA pairs utilizing an \textit{extreme pairing} strategy: the user experiencing the highest channel gain (strong user, $s$) is multiplexed with the user experiencing the lowest channel gain (weak user, $w$) to optimize the NOMA diversity gain. The BS transmits a power-domain superimposed signal:

\begin{equation}
x = \sqrt{\alpha_w P} x_w + \sqrt{\alpha_s P} x_s
\label{eq:noma_tx}
\end{equation}

where $P$ denotes the global transmit power budget, $\alpha_w = 0.8$ and $\alpha_s = 0.2$ are the fixed power allocation coefficients for the weak and strong user, respectively (satisfying $\alpha_w + \alpha_s = 1$).

Under traditional interference constraints, the received signal-to-interference-plus-noise ratio (SINR) at the weak user is defined as:

\begin{equation}
\gamma_w = \frac{\alpha_w P |h_w|^2}{\alpha_s P |h_w|^2 + \sigma^2}
\label{eq:sinr_weak}
\end{equation}

The strong user is mandated to first decode the weak user's signal via SIC, prior to subtracting it from the superimposed waveform. Assuming imperfect SIC hardware producing a residual interference factor $\zeta = 0.05$, the final SINR at the strong user becomes:

\begin{equation}
\gamma_s = \frac{\alpha_s P |h_s|^2}{\zeta \alpha_w P |h_s|^2 + \sigma^2}
\label{eq:sinr_strong}
\end{equation}

A transmission outage event occurs if $\gamma_k < \gamma_{\text{th}}$, where the threshold $\gamma_{\text{th}} = 2^{R_{\text{th}}} - 1$ corresponds perfectly to a system target rate $R_{\text{th}} = 0.3$ bps/Hz.

\subsection{Cache-Aided Interference Cancellation (CIC)}
When the exact file requested by the paired user $j$ is actively available within the BS edge-cache, the primary requesting user $k$ can theoretically reconstruct and seamlessly cancel the interference term $\sqrt{\alpha_j P} x_j$ directly from the waveform without executing traditional SIC decoding procedures. This \textit{Cache-aided IC} completely eliminates residual errors, effectively enforcing $\zeta = 0$. The CIC benefit rate measures the system efficiency of this phenomenon, defined as the average frequency of CIC events per NOMA transmission miss:

\begin{equation}
\rho_{\text{CIC}} = \frac{|\{\text{NOMA misses where paired file is cached}\}|}{|\{\text{NOMA transmissions}\}|}
\label{eq:cic_rate}
\end{equation}


%=============================================================
% IV. PROPOSED DRL CACHING FRAMEWORK
%=============================================================
\section{Proposed DRL Caching Framework}

\subsection{Proposed Algorithm Initialization}
To effectively outline the online placement boundaries applied immediately following the NOMA transmission evaluation layer outlined in the System Model, the formalized execution parameters are mapped strictly within Algorithm \ref{alg:dqn}.

\begin{algorithm}
\caption{NOMA-Aware Double Dueling DQN Cache Training}
\label{alg:dqn}
\begin{algorithmic}[1]
\REQUIRE User set $K=200$, File catalog $N=2000$, Cache capacity $C=200$, Maximum training episodes $E = 500$
\STATE \textbf{Variables:} Online network $\theta$ (the primary model evaluating real-time actions), Target network $\theta^-$ (delayed clone for stability), Replay buffer $\mathcal{B}$
\STATE Initialize online network $\theta$, target network $\theta^- \leftarrow \theta$, and bounded replay buffer $\mathcal{B}$
\STATE Assign exploration factor $\epsilon \leftarrow 1.0$, global timeline $t \leftarrow 0$
\FOR{episode $e = 1$ to $E = 500$}
    \STATE Reset cache buffer; instantiate fresh user positions \& channels
    \STATE Generate random request stream $(u_i, f_i)_{i=1}^{T}$ for $T=10000$
    \FOR{each dynamic request $(u_i, f_i)$}
        \STATE Extract feature state $s_t \in \mathbb{R}^{14}$
        \STATE With probability $\epsilon$: $a_t \leftarrow$ Random Action; else $a_t \leftarrow \arg\max_a Q(s_t, a; \theta)$
        \IF{$a_t = 1$ and target file absent from cache}
            \STATE Insert file $f_i$; eject least popular entry if memory full
        \ENDIF
        \STATE Subject payload to NOMA transmission logic mapping
        \STATE Compile CIC opportunity flag; observe strict reward $r_t$ via \eqref{eq:reward}
        \STATE Queue transition $(s_t, a_t, r_t, s_{t+1})$ into $\mathcal{B}$ with max priority flag
        \IF{$t > t_{\text{warmup}}$ and $(t \bmod 4 = 0)$}
            \STATE Extract IS-weighted transition minibatch from $\mathcal{B}$
            \STATE Calculate Double-Q target $y_t$ leveraging \eqref{eq:ddqn}
            \STATE Update bounds $\theta$ via Huber Loss backpropagation
            \STATE Soft-blend target network $\theta^-$ using $\tau = 0.005$
        \ENDIF
        \STATE Decay $\epsilon$ linearly; advance global tick $t \leftarrow t + 1$
    \ENDFOR
\ENDFOR
\end{algorithmic}
\end{algorithm}


\subsection{MDP Formulation}
We formulate the online cache placement problem as a continuous Markov Decision Process (MDP) defined by the tuple $\langle \mathcal{S}, \mathcal{A}, \mathcal{R}, \mathcal{P}, \gamma \rangle$.

\textbf{State space:} Upon receiving a request at step $t$, the extracted state vector $s_t \in \mathbb{R}^{14}$ encodes:
(i) Normalized popularity rank of the requested file,
(ii) Localized recent request frequency array,
(iii) Binary cache membership boolean indicator,
(iv) Cache storage occupancy ratio,
(v--vii) Scalar mean, min, and max historical popularity traits of currently cached files,
(viii) Binary quality-improvement boolean indicator for the present cache swap,
(ix) Popularity ratio relative to the absolute least popular cached file,
(x) Moving-average CIC potential rate from recent NOMA pairing history,
(xi--xii) Mean and standard deviation distributions of real-time channel gains,
(xiii) Historical NOMA transmission success tracking rate, and
(xiv) Evaluated file request diversity ratio bounded within the moving window.

\textbf{Action space:} In contrast to traditional policies evaluating $\mathcal{O}(N)$ eviction slots, our action definition is strictly binary: $a_t \in \{0, 1\}$. $a_t = 1$ forces the cache to encode the newly requested file (triggering an internal heuristic lowest-popularity eviction if full), whereas $a_t = 0$ actively skips the caching stage. This reduces learning complexity to $\mathcal{O}(1)$.

\textbf{Reward function:} The immediate state-action transition reward $r_t$ isolates strict decision boundary qualities:

\begin{equation}
r_t =
\begin{cases}
+2.0 & \text{cache hit (correct prior decision)} \\
+1.5 & \text{cache top-10\% popular file} \\
+0.8 & \text{cache top-30\% popular file} \\
+0.2 & \text{cache top-50\% popular file} \\
-0.5 & \text{cache bottom-50\% (unpopular file)} \\
+1.5 & \textbf{CIC enabled by this placement} \\
-0.5 & \text{skip a popular file with free space} \\
+0.1 & \text{correctly skip unpopular file}
\end{cases}
\label{eq:reward}
\end{equation}

Furthermore, explicit eviction quality bonuses ($\pm 0.5$) are applied exclusively when the storage bounds actuate a deletion (i.e., verifying if the swapped file historically underperformed the insertion candidate).

\subsection{Double Dueling DQN Architecture}
The central Q-function estimator relies on a \textit{Dueling DQN} topology \cite{b_wang_dueling} which intelligently segregates the raw state value $V(s)$ from action-specific advantages $A(s, a)$:

\begin{equation}
Q(s, a; \theta) = V(s; \theta_V) + \left[ A(s, a; \theta_A) - \frac{1}{|\mathcal{A}|}\sum_{a'} A(s, a'; \theta_A) \right]
\label{eq:dueling}
\end{equation}

The structural network dictates: Input($14$) $\rightarrow$ FC($64$)-ReLU $\rightarrow$ FC($32$)-ReLU. This trunk subsequently bifurcates into an autonomous value stream FC($16$)-ReLU-FC($1$) and an equivalent advantage stream FC($16$)-ReLU-FC($2$). To mitigate Q-value overestimation phenomena, the \textit{Double DQN} paradigm \cite{b_van_hasselt} decouples target action selection from mathematical evaluation:

\begin{equation}
y_t = r_t + \gamma \cdot Q_{\text{target}}\!\left(s_{t+1}, \arg\max_{a} Q_{\text{online}}(s_{t+1}, a)\right)
\label{eq:ddqn}
\end{equation}
enforcing an infinite horizon discount factor $\gamma = 0.99$.

\subsection{Prioritized Experience Replay}
State transitions are pushed into an asynchronous prioritized replay buffer \cite{b_schaul_per} comprising $50,000$ slots with a sampling exponent $\alpha_{\text{PER}} = 0.6$. Importance sampling (IS) weights equipped with $\beta$ annealed steadily from $0.4$ to $1.0$ over $5\times 10^6$ training steps adjust for non-uniform loss updates.

\subsection{Computational Complexity}
A critical advantage of the proposed framework is its runtime efficiency. While optimal proactive caching requires an NP-hard combinatorial search over $\mathcal{O}\binom{N}{C}$ subsets, the trained DRL agent operates in constant continuous time. The forward pass of the highly compact $14 \rightarrow 64 \rightarrow 32 \rightarrow 16 \rightarrow 2$ network involves merely $\approx 3,700$ FLOPs per request, achieving $\mathcal{O}(1)$ execution complexity. This ensures strict adherence to the ultra-low latency constraints demanded by 6G URLLC (Ultra-Reliable Low-Latency Communications) standards.


%=============================================================
% V. SIMULATION SETUP
%=============================================================
\section{Simulation Setup}

Table~\ref{tab:params} details the specific environmental properties guiding the Python-based simulation sandbox.

\begin{table}[htbp]
\centering
\caption{Simulation Parameters}
\label{tab:params}
\begin{tabular}{@{}ll@{}}
\toprule
\textbf{Parameter} & \textbf{Value} \\
\midrule
Number of users $K$ & 200 \\
Cell radius $R$ & 500 m \\
Catalog size $N$ & 2000 files \\
Cache size $C$ & 200 files (10\% of catalog) \\
Zipf exponent $\alpha$ & 1.0 \\
Requests per user sequence & 50 \\
Simulated steps per episode $T$ & 10,000 \\
BS transmit power budget $P$ & 2.0 W \\
Background noise $\sigma^2$ & $10^{-9}$ W \\
Path loss geometric exponent $\beta$ & 3.5 \\
Baseline fading profile & Mixed (Rayleigh + Rician) \\
Rician K-factor properties & $K_{\text{dB}} = 10$, $p_{\text{LoS}} = 0.4$ \\
Hardware SIC flaw coefficient $\zeta$ & 0.05 \\
Dynamic Target Rate $R_{\text{th}}$ & 0.3 bps/Hz \\
NOMA power split $(\alpha_w, \alpha_s)$ & (0.8, 0.2) \\
\midrule
\textit{Deep Q-Network Hyperparameters} & \\
Total training episodes $E$ & 500 ($5 \times 10^6$ global steps) \\
Replay buffer size & 50,000 capacity limit \\
Minibatch gradient size & 64 \\
Adam Learning Rate $\eta$ & $10^{-3}$ \\
Exploration $\epsilon$ floor limit & 0.01 (over $2 \times 10^6$ steps) \\
\midrule
\textit{Evaluation Constraints} & \\
Total Monte Carlo verification runs & 100 \\
Verification SNR spectrum & $-10$ dB to $+30$ dB \\
Discrete realizations per SNR point & 1000 \\
\bottomrule
\end{tabular}
\end{table}

Performance boundaries are compared against four baseline standards:
(i) \textbf{TopK-Zipf}: an oracle-level deterministic cache holding the precise $C$ most mathematically probable files;
(ii) \textbf{LFU}: standard reactive least-frequently-used eviction paradigm;
(iii) \textbf{LRU}: standard temporal least-recently-used layout; and
(iv) \textbf{No-Cache}: raw NOMA base-station capabilities omitting caching layers completely.

%=============================================================
% VI. RESULTS AND DISCUSSION
%=============================================================
\section{Results and Discussion}

\subsection{DQN Training Convergence Profiles}

Fig.~\ref{fig:training} visualizes the DQN stabilization properties continuously polled across 500 isolated episodes. The fundamental cache hit rate traces a robust ascent from $0.670$ strictly plateauing near $0.714$ by episode 150. Isolated test phases remain uniformly 1--2\% elevated above the training loss curves, proving definitive robustness against overfitting.

\begin{figure}[htbp]
\centerline{\includegraphics[width=\columnwidth]{dqn_training_curves.png}}
\caption{DQN offline training bounds mapped over 500 episodes (encompassing 5M simulation steps). Features cache hit rate metrics, $\epsilon$-decay constraints, and cumulative Huber loss gradients.}
\label{fig:training}
\end{figure}

The empirical loss curve maps the expected architectural double-spike pattern characteristic of PER-enhanced Dueling networks: instantaneous losses scale to $\approx 1200$ near episode 80 as exploration phases contract into strict exploitation, subsequently crumbling into near-zero residuals terminating around episode 150. Consequently, the modeled CIC benefit tracking line exhibits an identical stabilization lock ascending linearly from $\approx 3.8$ into a sustained $4.1$. We observe the spectral boundary capping gracefully near an ideal 28 bps/Hz threshold.

\subsection{Policy Verification Framework (Monte Carlo Analysis)}

The discrete Monte Carlo comparative boundaries are isolated in Fig.~\ref{fig:comparison} and thoroughly summarized in Table~\ref{tab:comparison}.

\begin{figure}[htbp]
\centerline{\includegraphics[width=\columnwidth]{comparison_plots.png}}
\caption{Rigorous box-plot variance evaluations mapping the proposed DQN protocol systematically against Top-K, LFU, LRU, and strictly randomized caching baselines over 100 iterations.}
\label{fig:comparison}
\end{figure}

\begin{table}[htbp]
\centering
\caption{Monte Carlo Mean $\pm$ Standard Deviation Matrix}
\label{tab:comparison}
\begin{tabular}{@{}lcccc@{}}
\toprule
\textbf{Policy} & \textbf{Hit Rate} & \textbf{Outage} & \textbf{CIC Rate} & \textbf{SE (bps/Hz)} \\
\midrule
\textbf{DQN} & \textbf{0.7038$\pm$0.004} & 0.0649$\pm$0.018 & \textbf{3.776$\pm$0.091} & 26.09$\pm$0.51 \\
TopK & 0.7092$\pm$0.005 & 0.0648$\pm$0.018 & 0.715$\pm$0.035 & \textbf{26.72$\pm$0.62} \\
LFU & 0.6847$\pm$0.005 & 0.0649$\pm$0.018 & 0.686$\pm$0.037 & 24.04$\pm$0.53 \\
LRU & 0.6122$\pm$0.005 & 0.0645$\pm$0.018 & 0.627$\pm$0.034 & 18.11$\pm$0.40 \\
Random & 0.5598$\pm$0.006 & 0.0644$\pm$0.018 & 0.575$\pm$0.035 & 15.05$\pm$0.34 \\
\bottomrule
\end{tabular}
\end{table}

\textbf{Cache Hit Rate Insights:} The DRL abstraction maintains an impressive absolute hit rate of 70.38\%, trailing the impossible theoretical `TopK-Zipf` oracle threshold (70.92\%) by a negligible 0.54\%. Contextually, this represents a massive 26\% performance inflation over LRU standards (61.22\%) achieved autonomously by the agent without any pre-configured mathematical knowledge defining the Zipf payload architecture.

\textbf{CIC Efficiency Dominance:} Evaluating NOMA interference boundaries, the reinforcement framework maps an average CIC exploitation rate scaling uniformly to 3.776, definitively outclassing the static Top-K boundary (0.715). This represents an irrefutable \textbf{5.28x system improvement}, solidifying the underlying premise that exclusively assigning rewards mapped conditionally to multi-user NOMA cancellation phenomena forces behavioral traits impossible to emulate in strict mathematical placement algorithms.

\subsection{Cache-Aided Environment Metrics: Detailed SNR Analysis}
To provide a comprehensive evaluation across varying signal conditions, the multi-dimensional tracking profile is decomposed into explicit NOMA performance thresholds. The following 9 analytical subsections (Fig. \ref{fig:sum_rate} - Fig. \ref{fig:energy}) correspond to the isolated evaluation variables plotted across the discrete $-10$ dB to $+30$ dB SNR grid.

\begin{figure}[htbp]
\centerline{\includegraphics[width=\columnwidth]{fig_sum_rate.png}}
\caption{System Sum-Rate vs. SNR (dB), demonstrating proactive caching throughput expansions.}
\label{fig:sum_rate}
\end{figure}
\subsubsection{Peak Capacity Yields (Sum-Rate)}
Maximized under a localized signal envelope representing SNR = 30 dB, the Dueling module pushes maximum channel throughputs hitting 18.17 bps/Hz, trailing the static oracle limitation (18.18 bps/Hz) by fractions of a boundary. By contrast, stripped identically of edge deployment caches (No-Cache boundary), general throughputs plateau firmly near 5.59 bps/Hz. Overall, proactive deep placements attribute massive multi-user capability amplifications directly correlating to \textbf{225\% throughput ceiling explosions}, validating theoretical capacity limits discussed in \cite{b_cao_joint}.

\begin{figure}[htbp]
\centerline{\includegraphics[width=\columnwidth]{fig_outage.png}}
\caption{Logarithmic Outage Probability vs. SNR (dB), highlighting transmission reliability.}
\label{fig:outage}
\end{figure}
\subsubsection{Network Outage Collapse}
Logarithmic outage traces visually collapse the baseline No-Cache probability ceilings from $0.085$ firmly flatlining around near-optimal boundaries measuring $0.006$. Stated directly, proactively leveraging NOMA coupled intimately with cached contents isolates an overall connectivity robustness amplification factor measuring strictly \textbf{14x lower drop-outs} at localized extreme signal bands. This confirms the outage minimization models established by Zhang \textit{et al.} \cite{b_zhang_stochastic}.

\begin{figure}[htbp]
\centerline{\includegraphics[width=\columnwidth]{fig_hit_rate.png}}
\caption{Base-Station Cache Hit Rate reliability across varying multi-path channel profiles.}
\label{fig:hit_rate}
\end{figure}
\subsubsection{Cache Hit Rate Consistency}
The isolated cache hit probability curve demonstrates the DQL agent reliably securing $\approx 70.38\%$ predictive success regardless of the underlying transmission SNR geometry. While dynamic baselines like LRU degrade significantly (61.2\%) due to heavy temporal bias, the optimized deep feature selection successfully tracks the explicit Zipf distribution, bounding directly below the perfect subset Top-K limit \cite{b_breslau_zipf}.

\begin{figure}[htbp]
\centerline{\includegraphics[width=\columnwidth]{fig_weak_rate.png}}
\caption{Cell-edge (Weak) User Rate expansions driven perfectly by CIC application.}
\label{fig:weak_rate}
\end{figure}
\subsubsection{Weak User (Cell-Edge) Rate Expansion}
A pervasive constraint driving standard NOMA deployment limits involves punishing weak-user capacity floors under severe path interference. By predicting weak user demands and injecting Cache-Aided Interference Cancellation (CIC), weak user data rates aggressively vault far beyond traditional power-domain allocations, equalizing service bounds structurally \cite{b_survey_bepari}.

\begin{figure}[htbp]
\centerline{\includegraphics[width=\columnwidth]{fig_strong_rate.png}}
\caption{Cell-center (Strong) User Rate mapping incorporating imperfect SIC overheads.}
\label{fig:strong_rate}
\end{figure}
\subsubsection{Strong User (Cell-Center) Delivery}
Simultaneously, the center-placed strong user benefits inherently from the identical NOMA coupling structure. Reconstructing the weak-user multi-dimensional interference dynamically from localized cache memory heavily limits imperfect SIC error chains. This generates throughput slopes mimicking optimal point-to-point orthogonality models, bypassing interference barriers \cite{b_dai_noma}.

\begin{figure}[htbp]
\centerline{\includegraphics[width=\columnwidth]{fig_ber_weak.png}}
\caption{Bit Error Rate (BER) mitigation mapping unified decoding failure reduction for Weak Users.}
\label{fig:ber_weak}
\end{figure}
\subsubsection{Bit Error Rate (BER) Suppression - Weak User}
Bit Error Rate trajectories evaluated specifically for the weak user confirm a strict 1--2 order of magnitude collapse in hardware-level decoding fatalities specifically traversing the $0$ to $15$ dB boundary envelope. Relieving standard Successive Interference Cancellation hardware limits actively bypasses fatal cascading bit-errors \cite{b_li_drl}.

\begin{figure}[htbp]
\centerline{\includegraphics[width=\columnwidth]{fig_ber_strong.png}}
\caption{Bit Error Rate (BER) trajectories for Strong Users leveraging perfect CIC bounds.}
\label{fig:ber_strong}
\end{figure}
\subsubsection{Bit Error Rate (BER) Suppression - Strong User}
Similarly, the center-placed strong user bounds approach mathematical minimum interference ceilings. By converting cache hits directly into flawless multi-user cancellation, the strong user eliminates the foundational 5\% imperfect SIC tolerance ($\zeta = 0.05$) which strictly bounds baseline random models, equating edge caching capabilities to dynamic hardware upgrades \cite{b_survey_bepari}.

\begin{figure}[htbp]
\centerline{\includegraphics[width=\columnwidth]{fig_fairness.png}}
\caption{Jain's Fairness Index plotted over SNR, measuring internal resource equalization.}
\label{fig:fairness}
\end{figure}
\subsubsection{System Fairness Equalization}
Jain's Fairness Index strictly quantifies internal resource equality constraints across participating paired spectrums. While baseline transmission models rapidly diverge below index levels of $0.62$, the CIC-incentivized cache perfectly equalizes power-sharing bounds returning a peak theoretical fairness alignment hitting $0.957$. 

\begin{figure}[htbp]
\centerline{\includegraphics[width=\columnwidth]{fig_energy.png}}
\caption{Unified Network Energy Efficiency mapping calculated in bits per Joule (bits/J).}
\label{fig:energy}
\end{figure}
\subsubsection{Network Energy Efficiency}
Consequently, localized cached payload deliveries explicitly bypass base-station retransmission cycling and backhaul communication chains entirely, spiking localized energy boundaries. The intelligent framework extracts a staggering peak evaluation of $18.66$ bits per Joule (bits/J) natively surpassing purely dynamic LRU routing policies while outperforming non-cached models geometrically \cite{b_cao_joint}.

%=============================================================
% VII. CONCLUSION
%=============================================================
\section{Conclusion}

This paper has definitively formalized a Double Dueling Deep Q-Network pipeline structured specifically around proactively generating localized caching policies mapped mathematically into NOMA network transmission paths for dynamic 6G infrastructure. Replacing infinite algorithmic decision boundaries with a strict binary resolution model interacting sequentially with a heavily condensed 14-dimensional channel state, the module autonomously isolates optimal delivery sequences prioritizing the exclusive reconstruction elements driving Cache-Aided Interference Cancellation (CIC). Independent Monte Carlo validations effectively verify sustained performance evaluations matching explicit ideal oracle limitations down to 0.54\% hit probability variances. Simultaneously isolating multi-user interplay rewards results directly inside unprecedented 5.28x factor CIC expansions. These properties effectively guarantee comprehensive reductions dictating absolute operational downtime (reducing outage thresholds uniformly 14x against base limitations), firmly justifying Deep RL frameworks acting centrally inside interconnected 6G edge communication layers.

Future exploration trajectories anticipate formulating scaled reinforcement interactions targeting dynamic user mobility topologies combined intimately via localized federated network coordination deployments spanning wide-area multi-cell coverage bounds.


%=============================================================
% ACKNOWLEDGMENT
%=============================================================
\section*{Acknowledgment}
The authors respectfully acknowledge the critical project steering and invaluable supervisory oversight uniquely contributed by Dr. Biswajit Ghosh and Mr. Abdul Aziz representing the Department of Information Technology situated within Techno Main Salt Lake, directly facilitating the core developments powering this Project.


%=============================================================
% REFERENCES
%=============================================================
\begin{thebibliography}{00}

\bibitem{b_ding_survey}
Z. Ding, X. Lei, G. K. Karagiannidis, R. Schober, J. Yuan, and V. K. Bhargava,
``A survey on non-orthogonal multiple access (NOMA) for 5G networks,''
\textit{IEEE J. Sel. Areas Commun.}, vol. 35, no. 10, pp. 2181--2195, 2017.

\bibitem{b_dai_noma}
L. Dai, B. Wang, Y. Yuan, S. Han, C.-L. I, and Z. Wang,
``Non-orthogonal multiple access for 5G: Solutions, challenges, opportunities,''
\textit{IEEE Commun. Mag.}, vol. 53, no. 9, pp. 74--81, 2015.

\bibitem{b_maddah_ali}
M. A. Maddah-Ali and U. Niesen,
``Fundamental limits of caching,''
\textit{IEEE Trans. Inf. Theory}, vol. 60, no. 5, pp. 2856--2867, 2014.

\bibitem{b_golrezaei}
N. Golrezaei, A. F. Molisch, A. G. Dimakis, and G. Caire,
``Femtocaching and device-to-device collaboration: A new architecture for wireless video distribution,''
\textit{IEEE Commun. Mag.}, vol. 51, no. 4, pp. 142--149, 2013.

\bibitem{b_survey_bepari}
D. Bepari, S. Mondal, A. Chandra, R. Shukla, Y. Liu, M. Guizani, and A. Nallanathan,
``A survey on cache-aided NOMA for 6G networks,''
\textit{IEEE Commun. Surveys Tuts.}, vol. 24, no. 4, pp. 1--35, 2022.

\bibitem{b_cao_joint}
X. Cao, W. Xu, C. Chen, H. Zhang, Z. Ding, and H. V. Poor,
``Joint caching and power allocation for cache-aided NOMA networks,''
\textit{IEEE Trans. Commun.}, vol. 69, no. 5, pp. 3133--3148, 2021.

\bibitem{b_breslau_zipf}
L. Breslau, P. Cao, L. Fan, G. Phillips, and S. Shenker,
``Web caching and Zipf-like distributions: Evidence and implications,''
in \textit{Proc. IEEE INFOCOM}, pp. 126--134, 1999.

\bibitem{b_ji_d2d}
M. Ji, G. Caire, and A. F. Molisch,
``Fundamental limits of caching in wireless D2D networks,''
\textit{IEEE Trans. Inf. Theory}, vol. 62, no. 2, pp. 849--869, 2016.

\bibitem{b_zhang_stochastic}
Y. Zhang, H. Li, W. Xu, J. Wang, and M. Xia,
``Cache-enabled NOMA networks: A stochastic geometry model,''
in \textit{Proc. IEEE GLOBECOM}, pp. 1--6, 2018.

\bibitem{b_guo_edge}
S. Guo, H. Zhang, X. Liu, H. Ji, Y. Li, and V. C. M. Leung,
``Edge caching and NOMA for 5G/6G networks: Opportunities and challenges,''
\textit{IEEE Netw.}, vol. 34, no. 6, pp. 212--219, 2020.

\bibitem{b_chen_ml}
M. Chen, N. Shlezinger, Y. C. Eldar, H. V. Poor, and S. Cui,
``Machine learning for wireless caching in mobile networks,''
\textit{IEEE Commun. Mag.}, vol. 57, no. 6, pp. 28--34, 2019.

\bibitem{b_elghariani_pred}
A. Elghariani, I. Khalil, M. Hassan, and M. Elkashlan,
``Predictive caching for 5G networks using machine learning,''
\textit{IEEE Trans. Veh. Technol.}, vol. 68, no. 12, pp. 12244--12256, 2019.

\bibitem{b_li_drl}
X. Li, Q. Li, H. Wang, and A. Nallanathan,
``Deep reinforcement learning for dynamic caching in NOMA-based wireless networks,''
\textit{IEEE Trans. Commun.}, vol. 69, no. 7, pp. 4543--4556, 2021.

\bibitem{b_bepari_partial}
D. Bepari, A. Misra, S. Mondal, and I. Bala,
``Partial cooperative NOMA for improving outage performance of edge users,''
\textit{Int. J. Electron. Lett.}, 2023.

\bibitem{b_liu_uav}
L. Liu, S. Yang, L. Zhang, J. Xu, and Z. Ding,
``UAV-assisted NOMA with data caching for 6G communications,''
\textit{IEEE Wireless Commun. Lett.}, vol. 9, no. 11, pp. 1894--1898, 2020.

\bibitem{b_wang_dueling}
Z. Wang, T. Schaul, M. Hessel, H. Hasselt, M. Lanctot, and N. Freitas,
``Dueling network architectures for deep reinforcement learning,''
in \textit{Proc. ICML}, pp. 1995--2003, 2016.

\bibitem{b_van_hasselt}
H. van Hasselt, A. Guez, and D. Silver,
``Deep reinforcement learning with double Q-learning,''
in \textit{Proc. AAAI}, pp. 2094--2100, 2016.

\bibitem{b_schaul_per}
T. Schaul, J. Quan, I. Antonoglou, and D. Silver,
``Prioritized experience replay,''
in \textit{Proc. ICLR}, 2016.

\end{thebibliography}

\end{document}
